\section{An order-three element in the 322-stem}
\label{sec:cube}

This section constructs the remaining positive class.  The construction is
independent of the height-six detector.  Its geometric input is the known
order-three Kervaire element in the 106-stem; its algebraic input is
Lemma~\ref{finite:main}, proved in Section~\ref{sec:finite}.  A truncated
extended cube first gives a map into a four-cell spectrum \(T\) containing
the Moore spectrum \(M=S/3\), with quotient \(\Sigma^4M\).  We construct
a map of Adams--Novikov filtration at least two with the same image in
that quotient; subtracting it gives a map that lifts to \(M\).  Keeping
track of that correction is essential both for the primary Steenrod
operation and for the additive order of the resulting sphere element.
Extended powers have long connected primary operations with unstable
and stable homotopy; see Toda~\cite{TodaExtended} and
May~\cite[Sections~2--3]{MayOperations}.
The construction below specifies the finite lens-space truncation
and the two Moore-spectrum corrections needed in this dimension.

\begin{theorem}\label{cube:order-three}
There is an element \(\theta\in\pi_{322}S_{(3)}\) detected by \(b_3\) in the
mod-three Adams spectral sequence such that \(3\theta=0\).
In particular, \(\theta\) has additive order exactly three.
\end{theorem}

\subsection{The strong criterion and the filtration conventions}

All spectra and maps in this section are localized at three.  Write \(S\)
for the localized sphere and let
\begin{equation}\label{cube:moore-triangle}
 S\xrightarrow{3}S\xrightarrow{i}M\xrightarrow{q}\Sigma S
\end{equation}
be the Moore cofibre sequence.  A power of a sphere class always denotes
its product in the stable homotopy ring of \(S\).

We use the following established form of the strong odd-primary Kervaire
criterion.  Our statement includes the change of indices and the duality
that will be needed below.

\begin{lemma}[Strong Kervaire criterion]\label{cube:strong-criterion}
Let \(j\geq0\) and \(m=3^{j+1}\).  The following conditions are equivalent:
\begin{enumerate}[label=\textup{(\roman*)}]
\item The detection coset of \(b_j\) contains an element of additive order
three in \(\pi_{4m-2}S\).
\item There is a stable map
\[
 \Sigma^{4m-2}M=P^{4m-1}(3)\longrightarrow S
\]
whose cone has nonzero reduced power \(\mathcal P^m\).
\item There is a stable map \(a:S^{4m-1}\to M\) for which
\[
 \mathcal P^m:H^0(C_a;\F_3)\longrightarrow H^{4m}(C_a;\F_3)
\]
is nonzero.
\end{enumerate}
In condition \textup{(iii)}, if \(qa=\Sigma\vartheta\), then
\(3\vartheta=0\) in actual sphere homotopy.  The strong criterion identifies
\(\vartheta\), up to the usual nonzero scalar normalization, as a nonzero
class detected by \(b_j\).
\end{lemma}

\begin{proof}
The equivalence of the first two conditions is Selick's strong
odd-primary criterion, in its stable Moore-space formulation
\cite[Theorems~1.1 and~3.3]{Selick}; see also the proof of
\cite[Theorem~1.1]{Amelotte}.  In the notation of that formulation the
stem is \(2m(p-1)-2\), and the detecting Adams class is \(b_{k-1}\) when
\(m=p^k\).  Thus at \(p=3\) and \(k=j+1\) the stem is \(4m-2\) and the
class is precisely \(b_j\).

For completeness, Spanier--Whitehead duality gives the last equivalence
without changing the primary test.  If
\(g:\Sigma^{4m-2}M\to S\), then \(DM\simeq\Sigma^{-1}M\) and
\[
 a=\Sigma^{4m-1}Dg:S^{4m-1}\longrightarrow M,\qquad
 C_a\simeq\Sigma^{4m}DC_g.
\]
The cohomology degrees of \(C_g\) and \(C_a\) are, respectively,
\(\{0,4m-1,4m\}\) and \(\{0,1,4m\}\).  Conjugation in the Steenrod
algebra sends \(\mathcal P^m\) to \(-\mathcal P^m\) modulo products of
positive reduced powers
\cite[Section~7, pp.~166--167]{Milnor}.  Such products vanish on the two cones,
because the necessary intermediate degrees are absent.  Consequently
duality preserves nonvanishing of the indicated operation.
Finally, \(3\vartheta=0\) follows directly from
\eqref{cube:moore-triangle}; the detection and nonvanishing are part of
the strong criterion.  Thus the criterion concerns an actual
order-three representative, rather than an order assertion about an
associated-graded group.
\end{proof}

Put
\[
 I=\operatorname{fib}(S\longrightarrow\BP),\qquad
 F_fZ=I^{\wedge f}\wedge Z,\qquad
 L_fZ=F_fZ/F_{f+1}Z.
\]
The Adams--Novikov filtration of a homotopy group is
\[
 F^f\pi_kZ=\im\bigl(\pi_kF_fZ\longrightarrow\pi_kZ\bigr).
\]
We retain the usual indices
\[
 E_r^{f,t}(Z)\Longrightarrow\pi_{t-f}Z,\qquad
 d_r:E_r^{f,t}(Z)\longrightarrow E_r^{f+r,t+r-1}(Z).
\]
In particular \(L_fS\simeq\BP\wedge I^{\wedge f}\).  Flatness of
\(\BP_*\BP\) identifies its homotopy groups with the desuspended
normalized cobar terms: they are torsion-free and supported in degrees
\(-f\bmod4\).  The groups for \(L_fM\) are their reductions modulo three.
It follows that
\begin{equation}\label{cube:sparseness}
 E_r^{f,t}(M)=0\quad\text{unless }t\equiv0\pmod4,
 \qquad E_2(M)=E_5(M),\quad E_6(M)=E_9(M).
\end{equation}
The tower description and its multiplicative pairings agree with the
usual Adams--Novikov spectral sequence; the unit-cube comparison and
the pairing conventions are recalled in Lemma~\ref{finite:pairing}.
The orientation \(\BP\to\HF\) gives the comparison with the ordinary
mod-three Adams tower.

These filtrations are finite in each fixed stem of a bounded-below
spectrum.  Indeed \(I\) is at least 3-connective: the unit is an
isomorphism on \(\pi_0\), and the localized sphere has no homotopy in
degrees one and two.  The connectivity of \(F_fZ\) therefore tends to
infinity with \(f\).  We shall use this finiteness when passing from
vanishing associated-graded groups to vanishing homotopy classes.

We also recall explicitly the elementary property of the first Adams
invariant that is used to control the corrections.

\begin{lemma}\label{cube:primary-invariant}
For a map \(u:A\to B\) inducing zero in mod-three homology, the short
exact sequence of Steenrod comodules
\[
 0\longrightarrow H_*B\longrightarrow H_*C_u
   \longrightarrow \Sigma H_*A\longrightarrow0
\]
is its ordinary Adams filtration-one invariant.  This invariant is
natural and additive in \(u\).  It vanishes for a composite of two maps
that induce zero in mod-three homology, and hence for a map of
Adams--Novikov filtration at least two.
\end{lemma}

\begin{proof}
The first Adams obstruction measures whether the displayed comodule
extension splits.  Cone maps give its naturality.  To add two maps,
first take their direct sum and then use the diagonal on the source
and the fold on the target; the corresponding pullback and pushout of
extensions give their Baer sum.  If \(u=vw\) and both \(v_*\) and
\(w_*\) vanish, the cone comparison from \(w\) to \(vw\) identifies
the extension of \(vw\) with the pushout of the extension of \(w\)
along \(v_*=0\).  This pushout is split.

The maps between two successive stages of the mod-three Adams tower
are homology-zero maps.  Thus a factorization through its second stage
has zero invariant.  The comparison \(\BP\to\HF\) supplies this
factorization for a map of Adams--Novikov filtration at least two.
\end{proof}

Whenever the cone has no ambiguity in the bottom and tested top
degrees, the coefficient of a primary reduced power between those
degrees is a component of this extension invariant.  In particular,
subtracting a map as in the last assertion of
Lemma~\ref{cube:primary-invariant} does not change that coefficient.

\subsection{A filtration-ten cube}

Set
\begin{equation}\label{cube:dimensions}
 t=108,\qquad n=t-1=107,\qquad N=3t-2=322.
\end{equation}
The classical Kervaire detection coset in \(\pi_{106}S\) contains
an element of order three detected by
\(b_2\); see \cite[Corollary~1.2 and its proof]{Amelotte} and
\cite[Table~A3.4]{Ravenel}.  By
Lemma~\ref{cube:strong-criterion}, choose
\[
 a_0:S^n\longrightarrow M
\]
with nonzero \(\mathcal P^{27}\) from the bottom to the top of its cone,
and write \(qa_0=\Sigma x\), where \(x\in\pi_{106}S\).

\begin{proposition}\label{cube:cubic-filtration}
The reduction of the cube satisfies
\[
 ix^3\in F^{10}\pi_{318}M.
\]
\end{proposition}

\begin{proof}
The class of \(a_0\) has Adams--Novikov filtration one.  Filtration
zero is excluded by the support of \(\pi_*\BP M\), while filtration
at least two is excluded by its nonzero primary test and
Lemma~\ref{cube:primary-invariant}.  Let
\[
 z\in E_2^{1,108}(M)
\]
be its permanent leading class.

We use the notation of Section~\ref{sec:finite}: \(v=v_1\),
\(a=[t_1]\), \(B\) is the reduced integral 3-Bockstein, and
\begin{align*}
 y&=\frac{\eta_R(v_2^9-v_1^8v_2^7)-(v_2^9-v_1^8v_2^7)}{v_1^9},\\
 t_7'&=\frac{\eta_R(v_2^7)-v_2^7}{v_1}
 \qquad\text{in the mod-three cobar complex}.
\end{align*}
By the unconditional basis calculation \eqref{finite:low-classes}, write
\[
 z=\lambda y+\nu t'_7+\mu v^{26}a,
 \qquad \lambda,\nu,\mu\in\F_3.
\]
After setting all positive coefficient generators \(v_i\) to zero,
these representatives specialize to \(\lambda t_1^{27}\).
Every one-cobar boundary also specializes to zero, since the
positive coefficient ideal is invariant.  The Thom comparison
sends \(t_1\), up to a unit, to the first reduced-power dual
generator; its conjugate Milnor convention changes only this
unit.  Hence \(\lambda\), up to a unit, is the
\(\mathcal P^{27}\) coefficient of \(C_{a_0}\), and is nonzero.
Rescale \(a_0\) to make \(\lambda=1\).
Lemma~\ref{finite:main} then gives
\begin{equation}\label{cube:leading-class}
 z=y+\mu v^{26}a,\quad\mu\in\F_3,\qquad B(y)^3=0.
\end{equation}

The boundary along \(q\) increases Adams--Novikov filtration by one
and induces the integral 3-Bockstein; this is the geometric boundary
theorem \cite[Theorem~2.3.4]{Ravenel}.  Here it can be seen directly
on two layers.  Put \(X=F_2S\), \(Y=F_1S\), and \(Z=Y/X\).
A lift of \(a_0\) to \(Y/3\) and an integral lift to \(Z\) of its
image in \(Z/3\) give a class in
\[
 Q=(Y/3)\mathbin{\times}_{Z/3}Z.
\]
The integral lift exists because the homotopy of \(Z\) is
torsion-free.  Comparing fibres over \(Z\) yields the triangle
\[
 X\xrightarrow{\,3\iota\,}Y
 \longrightarrow Q\longrightarrow\Sigma X,
\]
where \(\iota:X\to Y\) is the inclusion: the fibre of
\(Y\to Q\) identifies with the fibre of \(X\to X/3\),
and its map into \(Y\) is \(3\iota\).
Projection to \(Y/3\) identifies the resulting boundary with a lift
of \(x\) to \(F_2S\).  Projection to \(Z\) says that three times its
layer symbol is the \(d_1\) of the chosen integral lift.  Division
by three in the torsion-free layer gives exactly the stated
Bockstein, with a harmless convention-dependent sign.

After reducing again by \(i\), the filtration-two symbol is
\(\pm B(z)=\pm B(y)\).  Indeed the Bockstein is a derivation,
\[
 B(v)=\pm a,\qquad B(a)=0,\qquad a^2=0,
\]
so \(B(v^{26}a)=0\).  Multiplicativity of the sphere tower and its
pairing with the Moore tower puts \(ix^3\) in filtration at least
six, with filtration-six symbol represented by \(\pm B(y)^3\).
This symbol is zero by \eqref{cube:leading-class}.  Thus \(ix^3\)
has filtration greater than six.  Its stem is \(318\equiv2\bmod4\),
so \eqref{cube:sparseness} makes ten the next possible filtration.
\end{proof}

\subsection{The truncated extended cube}

Write \(\Sigma_3=C_3\rtimes C_2\), where the generator of \(C_2\)
acts on \(C_3\) by inversion.  Let \(\rho\) be the real permutation
representation on three letters and write
\[
 \rho=\mathbf1\oplus\lambda.
\]
Here \(\lambda\) is a real two-plane.  Restricted to \(C_3\) it is
the faithful complex character, while \(C_2\) acts by conjugation;
in particular its orientation character is \(\mathrm{sgn}\).

Let \(L^{(2)}=S^1\cup_3D^2\) be the two-skeleton of a lens-space
model for \(BC_3\).  A useful explicit cover is the subspace
\[
 \widetilde L^{(2)}
 =\{(z_1,z_2)\in S(\C^2):
       z_2\in\bigcup_{\zeta^3=1}\zeta\mathbf R_{\geq0}\}.
\]
Scalar multiplication by a third root of unity and complex
conjugation give the required \(C_3\rtimes C_2\)-action.  This
cover consists of three disks with their boundary circles
identified.  The \(C_3\)-quotient is \(L^{(2)}\), with its
degree-three attachment.

Take a free contractible \(C_2\)-space \(EC_2\), let \(C_3\) act
trivially on this extra factor, and put
\[
 \mathcal B=(L^{(2)})_{hC_2},\qquad
 \mathcal P=\widetilde L^{(2)}\times EC_2.
\]
Thus \(\mathcal P\to\mathcal B\) is a principal \(\Sigma_3\)-bundle.
A classifying map to \(B\Sigma_3\) gives a natural comparison
between
\begin{equation}\label{cube:truncated-functor}
 D^\circ(X)=\mathcal P_+\wedge_{\Sigma_3}X^{\wedge3},
 \qquad
 D_3(X)=E\Sigma_{3+}\wedge_{\Sigma_3}X^{\wedge3}.
\end{equation}
The permutation of smash factors is the usual permutation action
in spectra, including its signs on odd-dimensional spheres.

\begin{lemma}\label{cube:lens-thom}
The base \(\mathcal B\) has cohomological dimension at most two
for 3-local local coefficient systems, and
\[
 \widetilde H^*(\mathcal B;A)=0
\]
for every constant 3-local abelian group \(A\).
Moreover,
\begin{equation}\label{cube:thom-identifications}
 D^\circ(S^1)\simeq\Sigma^4M,\qquad
 D^\circ(S^n)\simeq\Sigma^NM.
\end{equation}
These identifications can be chosen so that
\begin{equation}\label{cube:quotient-cube}
 D^\circ(qa_0)=x^3 1_M:
       \Sigma^NM\longrightarrow\Sigma^4M.
\end{equation}
\end{lemma}

\begin{proof}
Since two is invertible, the spectral sequence for the homotopy
orbits by \(C_2\) takes invariants exactly.  The cellular complex
of \(L^{(2)}\) has dimension two, proving the assertion about
local coefficients.  With constant coefficients its possible
reduced cohomology groups are
\(\ker(3:A\to A)\) in degree one and \(A/3A\) in degree two.
Inversion acts by minus one on both, so neither has nonzero
\(C_2\)-invariants.  In particular the augmentation
\(\Sigma^\infty\mathcal B_+\to S\) is a 3-local equivalence:
its cofibre is bounded below and has zero integral localized
homology, hence is zero by the stable Hurewicz theorem.

The Thom spectrum of \(\rho\) over \(\mathcal B\) is
\(D^\circ(S^1)\).  Its orientation system is \(\mathrm{sgn}\).
On the lens skeleton, integral reduced homology is
\(\Z/3\) in degree one, and inversion acts by minus one.
Twisting by \(\mathrm{sgn}\) therefore retains that group
and removes the degree-zero group.  The Thom homology is
consequently just \(\Z/3\) in degree four.
To identify the spectrum, realize its lowest homotopy generator
by \(S^4\); it has order three by the lowest-degree Hurewicz
isomorphism and extends to \(\Sigma^4M\).
The extension is an integral homology equivalence and therefore
a 3-local equivalence of bounded-below spectra.

The difference bundle \((n-1)\rho\) is oriented, since \(n-1=106\)
is even, and is stably spherically trivial after localization at
three.  Indeed its spherical classifying map lifts to the
classifying space of degree-one stable sphere equivalences.
This space is simply connected and its second homotopy group
is \(\pi_1S_{(3)}=0\).  The subsequent Postnikov obstructions
vanish by the cohomological dimension of \(\mathcal B\).
There is no inverse-limit ambiguity: for a fixed homotopy
degree of the mapping space, the fibres of its Postnikov
tower have zero homotopy in that degree once their
Eilenberg--Mac Lane degree exceeds it by more than two.
Thus the required compatible nullhomotopies exist.
Trivializing this difference bundle identifies
\[
 D^\circ(S^n)\simeq
 \Sigma^{3(n-1)}D^\circ(S^1)
 \simeq\Sigma^{3n+1}M=\Sigma^NM.
\]

Before taking Thom spectra, \(D^\circ(qa_0)\) is a map
from the sphere bundle \(S^{n\rho}\) to \(S^\rho\).
After the difference bundle is trivialized, its fibrewise
mapping spectrum is the constant spectrum
\(S^{-3(n-1)}\): endomorphisms of the remaining invertible
sphere line identify canonically with scalars.
Sections of this constant spectrum are determined, on
homotopy classes, by evaluation at one point, because
\(\Sigma^\infty\mathcal B_+\simeq S\).
At that point the map is the ordinary smash cube of \(qa_0\),
namely the suspension of \(x^3\) in the indicated degrees.
This proves \eqref{cube:quotient-cube}, up to a 3-local unit
which is absorbed in the source identification.
\end{proof}

The same base calculation supplies the nonzero Euler pairing
that will detect a primary reduced power.  Namely
\begin{equation}\label{cube:euler-pairing}
 e(\lambda)\frown-:
 H_2(\mathcal B;\F_3^{\mathrm{sgn}})
 \longrightarrow H_0(\mathcal B;\F_3)
\end{equation}
is nonzero.  On restricting to the lens skeleton, \(e(\lambda)\)
is the mod-three first Chern class of the faithful \(C_3\)-line.
It is the nonzero Bockstein of a fundamental-group character.
Its evaluation on the mod-three two-cell is a unit.  The
orientation twist cancels the inversion sign, so taking
\(C_2\)-invariants preserves this pairing.

\begin{lemma}\label{cube:attachment}
There is a triangle
\begin{equation}\label{cube:T-triangle}
 \Sigma^3M\xrightarrow{g}M\xrightarrow{j_T}T
      \xrightarrow{\pi_T}\Sigma^4M
\end{equation}
and a map \(F:\Sigma^NM\to T\) such that
\[
 \pi_TF=x^3 1_M.
\]
Both \(j_T\) and the natural map \(D^\circ M\to T\) are
isomorphisms on bottom mod-three homology.
\end{lemma}

\begin{proof}
Filter \(D_3M\) by the number of upper Moore cells in its
three input factors.  The four successive layers, before
simplification, have stabilizers
\(\Sigma_{3-r}\times\Sigma_r\), \(0\leq r\leq3\).
The bottom layer is \(\Sigma^\infty B\Sigma_{3+}\).
For one upper cell, the \(\Sigma_2\)-stabilizer acts trivially
on its one-dimensional sphere; the layer is therefore
\(\Sigma S\) after localization at three.  For two upper
cells, interchanging the two odd spheres is the sign action.
Its \(\Sigma_2\)-homotopy orbits are zero at three, so this
layer disappears.  The top layer is \(D_3S^1\).

The attachment of the one-cell layer to the bottom is degree
three.  One can compute it in the cellular complex of
\(M^{\wedge3}\): a chosen one-upper-cell generator has
boundary three times the bottom generator.
Passing to \(\Sigma_3\)-coinvariants identifies the three
choices of this generator; it does not add their three
boundaries.  Thus the attachment has degree three, not nine.
Since the reduced suspension spectrum of the connected
space \(B\Sigma_3\) has zero \(\pi_0\), this attachment has
no reduced component.

It follows that the part \(K\) below the top layer has a
projection \(r:K\to M\) inducing an isomorphism on bottom
homology; explicitly, \(K\) is the cofibre of the bottom
degree-three map into \(\Sigma^\infty B\Sigma_{3+}\), and
the augmentation gives \(r\).
Pull back \(D_3M\to D_3S^1\) along
\(D^\circ S^1\to D_3S^1\), obtaining \(U\), and push out
its fibre along \(r\).  These constructions give a diagram
of triangles
\[
\begin{tikzcd}[column sep=large]
 K \arrow[r] \arrow[d,"r"'] &
 U \arrow[r] \arrow[d] &
 D^\circ S^1 \arrow[d,equal] \arrow[r] &
 \Sigma K \arrow[d,"\Sigma r"]\\
 M \arrow[r,"j_T"'] &
 T \arrow[r,"\pi_T"'] &
 D^\circ S^1 \arrow[r] &
 \Sigma M .
\end{tikzcd}
\]
Use \eqref{cube:thom-identifications} in the lower row
and desuspend its boundary to define \(g\).
The natural comparison square for \(D^\circ q\) gives
\(D^\circ M\to U\to T\); composing it with
\(D^\circ a_0\) gives \(F\).
Its quotient is \eqref{cube:quotient-cube}.
The bottom-homology claims follow because the comparisons
of base orbits preserve \(H_0\), the projection \(r\)
preserves the bottom generator, and the remaining quotient
\(\Sigma^4M\) has no homology near degree zero.
\end{proof}

\begin{proposition}\label{cube:primary-detection}
The cone of \(F:\Sigma^NM\to T\) has a nonzero operation
\[
 \mathcal P^{81}:H^0(C_F;\F_3)\longrightarrow
 H^{N+2}(C_F;\F_3).
\]
\end{proposition}

\begin{proof}
Throughout this proof homology has coefficients \(\F_3\).
We adapt Toda's diagonal cone construction
\cite[Section~1, Theorem~2]{TodaExtended} to the truncated functor
\(D^\circ\). Applying the diagonal cone nullhomotopy gives a natural map
\[
 C_{D^\circ a_0}\longrightarrow D^\circ(C_{a_0}).
\]
Compose it with the map \(D^\circ C_{a_0}\to D^\circ S^t\)
induced by \(C_{a_0}\to S^t\).  This composite is zero on
\(D^\circ M\), and therefore factors through
\(C_{D^\circ a_0}/D^\circ M\simeq\Sigma D^\circ S^n\).
The resulting map
\begin{equation}\label{cube:diagonal-thom}
 \Sigma D^\circ S^n\longrightarrow D^\circ S^t
\end{equation}
comes from the diagonal trivial line in the permutation
representation.  More explicitly, use the same interval
parameter in the three cone factors and then collapse
their ends.  Its normal representation is
\(\rho-\mathbf1=\lambda\).  The resulting Thom map is
therefore the Euler map for \(\lambda\).
Since \(N+2=3t=324\), the map on the tested top homology
is precisely the cap product \eqref{cube:euler-pairing};
it is nonzero.

The spectrum \(D^\circ M\) has homology only in degrees
at most five: its fibre \(M^{\wedge3}\) has homology in
degrees zero through three and its base has homological
dimension two.  Consequently the suspended top class
of \(D^\circ S^n\simeq\Sigma^NM\) gives a class
\(c\in H_{3t}(C_{D^\circ a_0})\).
Let \(u\in H_t(C_{a_0})\) be the top generator.
In the homology spectral sequence over \(\mathcal B\),
the fibre of \(D^\circ C_{a_0}\) has tensor-product
homology with degrees chosen from \(\{0,1,t\}\).
In total degree \(3t\), the sole possible term is
\[
 H_0\bigl(\mathcal B;\F_3\{u^{\otimes3}\}\bigr);
\]
the action is trivial since \(t\) is even.
There is no possible differential into or out of this
term, nor another term in its total degree.
Projection to \(D^\circ S^t\) identifies it with the
fibre generator there.  By \eqref{cube:diagonal-thom},
the image of \(c\) is a nonzero multiple of
\(u^{\otimes3}\).

The bottom fibre class survives as well.  The Cartan formula
for reduced powers \cite[Section~2, Property~2.5]{Milnor} on the
fibre gives
\[
 \mathcal P_*^{81}(u^{\otimes3})
       =(\mathcal P_*^{27}u)^{\otimes3}\ne0.
\]
Indeed, to reach degree zero each tensor factor must
be lowered from degree \(t=108\) by \(\mathcal P_*^{27}\);
all other Cartan terms have zero contribution to the
bottom.  Naturality shows that
\(\mathcal P_*^{81}c\) is nonzero in
\(H_0(C_{D^\circ a_0})\).

Finally, the map \(D^\circ M\to T\) and the identity on
the source induce \(C_{D^\circ a_0}\to C_F\).
It is an isomorphism on the bottom and on the tested top
homology.  Applying naturality again proves the assertion.
\end{proof}

\subsection{Endomorphisms and the two obstruction components}

\begin{lemma}\label{cube:endomorphisms}
There is a unital multiplication \(\mu:M\wedge M\to M\),
and evaluation on the bottom cell admits a splitting
\begin{equation}\label{cube:end-splitting}
 \mathcal E=F(M,M)\simeq M\vee\Sigma^{-1}M.
\end{equation}
Write \(\varepsilon:\mathcal E\to M\) for evaluation,
\(s:M\to\mathcal E\) for its section, and
\(\tau:\mathcal E\to\Sigma^{-1}M\) for the other projection.
The second summand includes by precomposition with
\(q:M\to\Sigma S\), and \(s(i)=1_M\).
No associativity of \(\mu\) is required.
\end{lemma}

\begin{proof}
The Moore sequence gives \(\pi_1M=0\) and
\(\pi_0M=\F_3\).  Applying \([-,M]\) to that sequence
shows that restriction to the bottom identifies
\([M,M]\) with \(\pi_0M\).  In particular \(3\,1_M=0\).
The cofibre sequence
\[
 M\xrightarrow{3}M\longrightarrow M\wedge M
\]
then allows the identity of \(M\) to extend to a
multiplication \(\mu\) with one unit law.  The other
unit composite has the same bottom restriction as
the identity, so it is the identity by the same
restriction calculation.

Applying \(F(-,M)\) to \eqref{cube:moore-triangle}
identifies the fibre of
\(\varepsilon:F(M,M)\to F(S,M)=M\) with
\(F(\Sigma S,M)=\Sigma^{-1}M\).
Adjunction applied to \(\mu\) gives the section \(s\).
Its value on \(i\) is the identity by the unit law.
Splitting this fibre sequence proves the assertions.
\end{proof}

The correction will be made inside \(F_2T\).  It is
important to distinguish its tower from the full
Adams--Novikov tower.

\begin{lemma}\label{cube:relative-tower}
Let \(\widetilde E_r\) be the spectral sequence of
\[
 \cdots\longrightarrow F_4M\longrightarrow F_3M
       \longrightarrow F_2M,
\]
with absolute filtration indices retained.  Its
comparison \(c_r:\widetilde E_r^{f,t}\to E_r^{f,t}(M)\)
has the following properties:
\[
\begin{array}{c|c|c}
\text{page}&\text{surjective}&\text{isomorphic}\\ \hline
5&f\geq2&f>2\\
6=9&f\geq2&f\geq7\\
10&f\geq2&f\geq11.
\end{array}
\]
These assertions hold also for maps induced by fixed
endomorphisms of \(M\), with their appropriate degree
shifts.
\end{lemma}

\begin{proof}
The layer complex agrees with the full \(E_1\)-complex
above filtration two.  In filtration two it has the
same cycles but omits incoming \(d_1\)-boundaries.
This gives the stated comparison on \(E_2\), hence on
\(E_5\) by sparseness.

For clarity, a full \(d_r\)-cycle lifts to a cycle
whenever comparison is surjective in its slot and
injective in its outgoing slot: the differential of
any lift maps to zero and must therefore vanish.
Comparison on the resulting homology is injective
as well if it is injective in the current slot and
surjective in the incoming slot, because a full
boundary can then be lifted.

Apply this first to \(d_5\).  An outgoing slot from
filtration \(f\geq2\) has filtration at least seven,
where \(c_5\) is an isomorphism.  Thus \(c_6\) is
surjective for every \(f\geq2\).  If \(f\geq7\), the
incoming slot has filtration \(f-5\geq2\), so the
same argument gives an isomorphism.  Sparseness
identifies this page with \(E_9\).
For \(d_9\), all outgoing slots again compare
isomorphically.  When \(f\geq11\), the incoming
filtration \(f-9\) is at least two, proving the
last row.  Everything is natural in maps of the
underlying tower.
\end{proof}

Write the filtration of the abutment of this restricted
tower as
\[
 \mathcal F^f\pi_kF_2M
   =\im\bigl(\pi_kF_fM\longrightarrow\pi_kF_2M\bigr),
 \qquad f\geq2.
\]
It is finite by connectivity, just as the full filtration.

\begin{proposition}\label{cube:correction}
There is a map \(F':\Sigma^NM\to F_2T\) whose composite
to \(T\) has the same quotient under \(\pi_T\) as \(F\).
\end{proposition}

\begin{proof}
Choose a lift
\(\widehat c\in\pi_{318}F_{10}M\) of \(ix^3\) using
Proposition~\ref{cube:cubic-filtration}, and let \(c\)
be its image in \(\pi_{318}F_2M\).  Smash the section
in Lemma~\ref{cube:endomorphisms} with \(I^{\wedge2}\)
and put
\[
 w=s_*(c)\in\pi_{318}F_2\mathcal E.
\]
After forgetting filtration, \(w\) is exactly
\(x^3 1_M\), because \(s(i)=1_M\).

Postcomposition with \(g:\Sigma^3M\to M\) defines
\[
 \gamma:\Sigma^3\mathcal E\longrightarrow\mathcal E.
\]
The obstruction to lifting \(w\) through the triangle
\eqref{cube:T-triangle} smashed with \(I^{\wedge2}\)
is
\[
 \gamma_*(w)\in\pi_{321}F_2\mathcal E.
\]
We show separately that both components of this
obstruction under \eqref{cube:end-splitting} vanish.

\emph{The first component.}
Forgetting filtration makes the whole obstruction
zero, since \(F\) is already an actual lift of
\(x^3 1_M\).  Its first component lies in
\(\pi_{321}F_2M\), and the map
\[
 \pi_{321}F_2M\longrightarrow\pi_{321}M
\]
is injective.  To see this, the quotient \(M/F_2M\)
has its two layers \(L_0M\) and \(L_1M\), both with
zero \(\pi_{322}\): their support degrees are,
respectively, \(0\) and \(3\) modulo four, whereas
\(322\equiv2\).  Thus \(\pi_{322}(M/F_2M)=0\), and
the long exact sequence gives the injection.
The first component is therefore zero already in
the restricted tower.

\emph{The second component.}
The fixed composite
\[
 \Sigma^3M\xrightarrow{\Sigma^3s}\Sigma^3\mathcal E
     \xrightarrow{\gamma}\mathcal E
     \xrightarrow{\tau}\Sigma^{-1}M
\]
suspends to a map \(h:\Sigma^4M\to M\).
Under the splitting, the second obstruction is
precisely \(h_*(c)\in\pi_{322}F_2M\).
The map \(h\) on \(\BP_*\)-homology is multiplication
by a scalar multiple of \(v_1\), since a
\(\BP_*\)-linear degree-four map
\(\BP_*M\to\BP_*M\) is determined by its value in
\(\BP_4M=\F_3\{v_1\}\).
This is a comodule map: \(v_1\) is primitive modulo three,
by \(\eta_R(v_1)=v_1+3t_1\).
Functoriality of the normalized cobar complex therefore
identifies \(E_2(h)\) with that literal coefficient
multiplication.  The map of Adams towers passes this
identification to \(E_5\) and \(E_6\); the coefficient
\(v_1\) is permanent by the target-vanishing argument in
Section~\ref{sec:finite}.
No multiplicativity assumption on \(h\) is used.
Its action on the full \(E_6\)-page in the
relevant slots is consequently a scalar multiple of
\begin{equation}\label{cube:v-obstruction}
 v_1:E_6^{10,328}(M)\longrightarrow E_6^{10,332}(M),
\end{equation}
which is zero by Lemma~\ref{finite:main}.
Lemma~\ref{cube:relative-tower} is an isomorphism on
both these \(E_6\)-slots.  Consequently \(h\) induces
zero on the filtration-ten associated-graded map
of the restricted abutment.  Since
\(c\in\mathcal F^{10}\pi_{318}F_2M\), this proves
\[
 h_*(c)\in\mathcal F^{11}\pi_{322}F_2M
          =\mathcal F^{14}\pi_{322}F_2M.
\]
The equality follows from sparseness: in stem 322
only filtrations congruent to two modulo four can
occur.

For every such filtration \(f\geq14\),
Lemma~\ref{finite:main}, including its high-filtration
estimate, gives
\[
 E_{10}^{f,322+f}(M)=0.
\]
The comparison of Lemma~\ref{cube:relative-tower}
is an isomorphism there, so the restricted
\(E_{10}\)-groups vanish as well.  Every
associated-graded quotient of
\(\mathcal F^{14}\pi_{322}F_2M\) is therefore zero.
Increasing connectivity makes this filtration
finite, and hence
\(\mathcal F^{14}\pi_{322}F_2M=0\).
This proves that the second component vanishes.

Notice why the two page arguments were separated:
an incoming full \(d_9\) at filtration ten may
start in filtration one, which the restricted
tower omits.  No \(E_{10}\)-comparison in
filtration ten was asserted or used.  The zero
map \eqref{cube:v-obstruction} removes that
filtration on \(E_6\) before the higher
\(E_{10}\)-vanishing is applied.

We have now proved \(\gamma_*(w)=0\).
The triangle \eqref{cube:T-triangle} smashed with
\(I^{\wedge2}\) supplies the desired lift.
More explicitly, \(M\) is dualizable, so
\[
 F(M,I^{\wedge2}\wedge Z)
   \simeq I^{\wedge2}\wedge F(M,Z).
 \]
The exact sequence obtained by applying \(F(M,-)\)
to that triangle therefore identifies the
vanishing obstruction just calculated with the
existence of \(F':\Sigma^NM\to F_2T\).
Its quotient represents \(w\), whose underlying
map is \(x^3 1_M=\pi_TF\).
\end{proof}

\subsection{The sphere element and its order}

\begin{proof}[Proof of Theorem~\ref{cube:order-three}]
Let \(\overline F'\) be the composite of the map in
Proposition~\ref{cube:correction} with \(F_2T\to T\).
Its quotient equals that of \(F\), so the difference
has a lift
\[
 Z:\Sigma^NM\longrightarrow M,\qquad
 j_TZ=F-\overline F'.
\]
The map \(\overline F'\) has Adams--Novikov filtration
at least two.  The maps \(F\), \(\overline F'\), and
\(F-\overline F'\) induce zero on mod-three homology for degree reasons, so
Lemma~\ref{cube:primary-invariant} and additivity
show that the primary invariant of
\(F-\overline F'\) equals that of \(F\).
Since \(j_T\) is an isomorphism on bottom homology,
the cone comparison from \(Z\) to \(j_TZ\) is an
isomorphism on the bottom and on the tested top
degrees.  Proposition~\ref{cube:primary-detection}
thus gives a nonzero
\[
 \mathcal P^{81}:H^0(C_Z;\F_3)\longrightarrow
                      H^{324}(C_Z;\F_3).
\]

Use the splitting of Lemma~\ref{cube:endomorphisms}
to write the class of \(Z\) as the sum of its
section component and its quotient-precomposition
component.  The section component is \(s_*(z_0)\)
for some \(z_0\in\pi_NM=\pi_{322}M\).
This class has Adams--Novikov filtration at
least two: the degree-322 homotopy of both
\(L_0M\) and \(L_1M\) is zero, so it lifts
successively through \(F_1M\) and \(F_2M\).
Applying the section to that lift shows that
the corresponding map \(\Sigma^NM\to M\)
also factors through \(F_2M\).
It can therefore be subtracted without
altering the primary test.

The remaining component factors through the
upper quotient of its source:
\[
 \Sigma^NM\xrightarrow{\Sigma^Nq}S^{N+1}
       \xrightarrow{a_1}M.
\]
The cone comparison induced by \(\Sigma^Nq\)
preserves the bottom class and the top class
in degree \(N+2=324\).  Hence \(C_{a_1}\)
has nonzero \(\mathcal P^{81}\) from its
bottom to its top.  Apply
Lemma~\ref{cube:strong-criterion} with
\(m=81=3^4\): the map
\(a_1:S^{323}\to M\) gives an element
\(\theta\in\pi_{322}S\) detected by \(b_3\).
The actual Moore boundary satisfies
\[
 qa_1=\Sigma\theta,\qquad 3\theta=0.
\]
Nonzero detection implies \(\theta\ne0\);
therefore its additive order is exactly three.
\end{proof}

\begin{remark}
The map into \(T\) and its nonzero primary operation
are supplied directly by the extended cube.
The finite Adams--Novikov lemma is used to produce
the correction inside \(F_2T\), thereby obtaining
the map into \(M\).  This is the step that determines
the additive order in the sphere, and it does not
follow merely from a differential-free
associated-graded symbol.
\end{remark}
